% =========================================================================
% Figura: Representación de moléculas en SMILES (Aspirina 2D vs. SMILES)
% Paleta de colores accesible Okabe-Ito definida en Configurations/01-Colors.sty
% =========================================================================

\begin{figure}[htbp]
\centering
\resizebox{0.95\linewidth}{!}{%
\begin{tikzpicture}[
    font=\sffamily,
    >=Stealth,
    x=1cm, y=1cm
]

% =========================================================================
% CONTENEDOR PRINCIPAL
% =========================================================================
\draw[draw=cardBorder, fill=white, rounded corners=10pt, line width=0.9pt] 
    (-0.5, -0.2) rectangle (15.8, 6.9);

% =========================================================================
% PANEL IZQUIERDO: ESTRUCTURA 2D (GRAFO MOLECULAR)
% =========================================================================
\node[anchor=north west, font=\bfseries\small, text=darkText] at (-0.1, 6.55) {Grafo 2D};
\node[anchor=north west, font=\scriptsize, text=grayText] at (-0.1, 6.15) {Molécula de aspirina};

% Coordenadas del anillo bencénico (Centro en (2.2, 2.9), R = 1.18 cm)
\coordinate (C1) at (3.22, 3.49); % C1 (ramificación 1)
\coordinate (C6) at (2.20, 4.08);
\coordinate (C5) at (1.18, 3.49);
\coordinate (C4) at (1.18, 2.31);
\coordinate (C3) at (2.20, 1.72);
\coordinate (C2) at (3.22, 2.31); % C2 (ramificación 2)

% Coordenadas ramificación lateral 1 (arriba a la derecha)
\coordinate (O_est)  at (3.98, 4.18);
\coordinate (C_act)  at (4.80, 4.85);
\coordinate (O_carb) at (4.80, 5.92);
\coordinate (C_met)  at (5.85, 4.28);

% Coordenadas ramificación lateral 2 (abajo a la derecha)
\coordinate (C_acd)  at (4.15, 1.76);
\coordinate (O_acd)  at (5.08, 2.30);
\coordinate (OH_acd) at (4.15, 0.62);

% --- HALOS COLOREADOS DETRÁS DE CADA FRAGMENTO ---
% 1. Halo Ramificación lateral (Naranja)
\filldraw[fill=figOrangeTint, draw=figOrange, line width=1.2pt, dashed, rounded corners=8pt]
    (3.55, 3.65) -- (3.55, 6.30) -- (6.45, 6.30) -- (6.45, 3.65) -- cycle;

% 2. Halo Ciclo del grafo (Azul)
\filldraw[fill=figBlueTint, draw=figBlue, line width=1.2pt, dashed, rounded corners=8pt]
    (0.70, 1.25) -- (0.70, 4.55) -- (3.48, 4.55) -- (3.48, 1.25) -- cycle;

% 3. Halo Ramificación terminal (Verde)
\filldraw[fill=figGreenTint, draw=figGreen, line width=1.2pt, dashed, rounded corners=8pt]
    (3.65, 0.15) -- (3.65, 2.55) -- (5.60, 2.55) -- (5.60, 0.15) -- cycle;

% --- ENLACES Y ÁTOMOS 2D ---
% Anillo (Kekulé con enlaces dobles alternados)
\draw[line width=1.5pt, draw=darkText] (C1) -- (C6) -- (C5) -- (C4) -- (C3) -- (C2) -- cycle;
\draw[line width=1.1pt, draw=darkText] ($(C2)!0.15!(C1) + (-0.14, 0)$) -- ($(C1)!0.15!(C2) + (-0.14, 0)$);
\draw[line width=1.1pt, draw=darkText] ($(C5)!0.15!(C6) + (0.07, -0.12)$) -- ($(C6)!0.15!(C5) + (0.07, -0.12)$);
\draw[line width=1.1pt, draw=darkText] ($(C3)!0.15!(C4) + (-0.07, 0.12)$) -- ($(C4)!0.15!(C3) + (-0.07, 0.12)$);

% Índices '1' de cierre de ciclo
\node[font=\fontsize{6.5}{7.5}\selectfont\bfseries, text=white, circle, fill=figBlue, inner sep=1pt] at ($(C1)+(-0.35, 0)$) {1};
\node[font=\fontsize{6.5}{7.5}\selectfont\bfseries, text=white, circle, fill=figBlue, inner sep=1pt] at ($(C2)+(-0.35, 0)$) {1};

% Ramificación lateral 1
\draw[line width=1.5pt, draw=darkText] (C1) -- (O_est) -- (C_act);
\draw[line width=1.3pt, draw=darkText] ($(C_act)+(-0.06,0)$) -- ($(O_carb)+(-0.06,0)$);
\draw[line width=1.3pt, draw=darkText] ($(C_act)+(0.06,0)$) -- ($(O_carb)+(0.06,0)$);
\draw[line width=1.5pt, draw=darkText] (C_act) -- (C_met);

\node[font=\footnotesize\bfseries, text=darkText, fill=figOrangeTint, inner sep=1.5pt, rounded corners=1pt] at (O_est) {O};
\node[font=\footnotesize\bfseries, text=darkText, fill=figOrangeTint, inner sep=1.5pt, rounded corners=1pt] at (O_carb) {O};
\node[font=\footnotesize\bfseries, text=darkText, fill=figOrangeTint, inner sep=1.5pt, rounded corners=1pt] at (C_met) {CH$_3$};

% Ramificación terminal 2
\draw[line width=1.5pt, draw=darkText] (C2) -- (C_acd);
\draw[line width=1.3pt, draw=darkText] ($(C_acd)+(0.05, 0.07)$) -- ($(O_acd)+(0.05, 0.07)$);
\draw[line width=1.3pt, draw=darkText] ($(C_acd)+(-0.05, -0.07)$) -- ($(O_acd)+(-0.05, -0.07)$);
\draw[line width=1.5pt, draw=darkText] (C_acd) -- (OH_acd);

\node[font=\footnotesize\bfseries, text=darkText, fill=figGreenTint, inner sep=1.5pt, rounded corners=1pt] at (O_acd) {O};
\node[font=\footnotesize\bfseries, text=darkText, fill=figGreenTint, inner sep=1.5pt, rounded corners=1pt] at (OH_acd) {OH};

% =========================================================================
% PANEL DERECHO: STRING SMILES (SERIALIZACIÓN LINEAL)
% =========================================================================
\node[anchor=north west, font=\bfseries\small, text=darkText] at (8.5, 6.55) {Secuencia SMILES};
\node[anchor=north west, font=\scriptsize, text=grayText] at (8.5, 6.15) {Serialización lineal del grafo};

% Cápsula SMILES completa destacada
\node[
    draw=cardBorder,
    fill=cardBorder!15,
    rounded corners=6pt,
    line width=0.8pt,
    inner sep=6pt,
    anchor=north west
] (smiles_bar) at (8.5, 5.65) {
    {\fontsize{14}{17}\selectfont\ttfamily
    \colorbox{figOrangeTint}{\textbf{\color{figOrange}CC(=O)O}}%
    \colorbox{figBlueTint}{\textbf{\color{figBlue}c1ccccc1}}%
    \colorbox{figGreenTint}{\textbf{\color{figGreen}C(=O)O}}%
    }
};

% 3 Bloques visuales correspondientes a cada fragmento del grafo
% Bloque 1: Ramificación lateral (Naranja)
\node[
    draw=figOrange,
    fill=figOrangeTint,
    rounded corners=5pt,
    line width=1.1pt,
    anchor=north west,
    minimum width=6.5cm,
    inner sep=5pt
] (box_est) at (8.5, 4.30) {
    \makebox[2.4cm][l]{{\fontsize{12}{14}\selectfont\ttfamily\textbf{\color{figOrange}CC(=O)O}}}%
    {\small\textbf{\color{darkText}Ramificación lateral}}
};

% Bloque 2: Estructura cíclica (Azul)
\node[
    draw=figBlue,
    fill=figBlueTint,
    rounded corners=5pt,
    line width=1.1pt,
    anchor=north west,
    minimum width=6.5cm,
    inner sep=5pt
] (box_rng) at (8.5, 2.95) {
    \makebox[2.4cm][l]{{\fontsize{12}{14}\selectfont\ttfamily\textbf{\color{figBlue}c1ccccc1}}}%
    {\small\textbf{\color{darkText}Estructura cíclica} (\texttt{1...1})}
};

% Bloque 3: Ramificación terminal (Verde)
\node[
    draw=figGreen,
    fill=figGreenTint,
    rounded corners=5pt,
    line width=1.1pt,
    anchor=north west,
    minimum width=6.5cm,
    inner sep=5pt
] (box_acd) at (8.5, 1.60) {
    \makebox[2.4cm][l]{{\fontsize{12}{14}\selectfont\ttfamily\textbf{\color{figGreen}C(=O)O}}}%
    {\small\textbf{\color{darkText}Ramificación terminal}}
};

% =========================================================================
% FLECHAS DE CONEXIÓN ENTRE FRAGMENTO 2D Y BLOQUE SMILES
% =========================================================================
\draw[figOrange, line width=1.5pt, -{Stealth[length=6pt, width=4.5pt]}]
    (6.45, 5.0) to[out=5, in=175] (box_est.west);

\draw[figBlue, line width=1.5pt, -{Stealth[length=6pt, width=4.5pt]}]
    (3.48, 3.15) to[out=10, in=175] (box_rng.west);

\draw[figGreen, line width=1.5pt, -{Stealth[length=6pt, width=4.5pt]}]
    (5.60, 1.45) to[out=-5, in=185] (box_acd.west);

\end{tikzpicture}%
}
\caption{Representación de la molécula de aspirina en notación SMILES y correspondencia con su grafo 2D.}
\label{fig:smiles_aspirina}
\end{figure}
